\section{Original logical rank and the normal-type checking variant}
\label{sec:p07-ch88-checking}

This section uses the corrected indexed binding and the full/restricted
source rules of \cref{sec:p07-ch88}. The original checking discussion is
on pp.~104--105; the separate no-conversion, normal-type presentation is
\S6.1, pp.~111--112 \cite{CoquandHuet1988}. The latter has a complete
checker below, with a direct adequacy proof for its declared domain.
It is not silently identified with the full calculus on unchanged raw
terms. All additional written rules are on
Sheet~\ref{sheet:p07-ch88-normal-types}.

Named expressions in the examples denote the indexed inverse of the
role-sensitive encoding \(T\), with a fixed external variable vector.
This is the proved raw binding correspondence of
\cref{thm:p07-ch88-binding-transport}, not an assumed inverse of target
typing. In particular the named example variables are actual bound or
external references, not the dummy letters of indexed binder labels.

\begin{hblemma}{thm:p07-ch88-check-generation}{proved=annex:p7.ch88.check-generation}{Source context regularity and conversion-stripping generation}
Every derived validity, typing or typed-equality judgment in either
corrected source variant has a valid environment. Its declarations have
their actual older-prefix certificates: a logical-context annotation is
valid there, or an object annotation has type \(\ast\) there.

If \(\Gamma\vdash M:P\), stripping final typing conversions exposes
the common rule for the literal subject constructor. A variable's initial
type is its lookup; an object product has initial type \(\ast\) and a
body premise typed \(\ast\); a lambda has a body typing and the
corresponding literal product type; an application has function/argument
premises and initial type \([N/x]B\). The initial type is joined to
\(P\) by the actual finite chain of source typing-conversion premises,
allowing a zero-length chain. This statement does not assume that every
raw equality endpoint is itself a sorted type.
\end{hblemma}

\begin{hbtheorem}{thm:p07-ch88-check-substitution}{proved=annex:p7.ch88.check-substitution}{Actual source thinning and dependent substitution}
In either corrected ordinary source variant, a judgment may be thinned
to a valid environment containing its declarations in the same relative
order, with the same annotations and fresh additional declarations.
The indexed conclusion uses the corresponding reference relocation.
If \(\Gamma[x:A];\Theta\vdash J\) and \(\Gamma\vdash N:A\),
where \(N\) is a source object, then
\(\Gamma;([N/x]\Theta)\vdash[N/x]J\). This covers validity, typing
and typed equality; every suffix annotation is substituted in its own
older prefix. No beta subject-reduction or inverse CC typing premise
is needed.
\end{hbtheorem}

\begin{hblemma}{thm:p07-ch88-check-rank}{proved=annex:p7.ch88.check-rank}{The source rank-invariance lemma}
For the source rank \(\delta\) of
\cref{def:p07-ch88-logical-reduction}, an object-argument substitution
preserves the rank of every logical-context term. Source typed equality
preserves the object/logical-context class and \(\delta\), in either
corrected variant. This supplies the actual source Lemma~6's conclusion,
but not the subsequent claim about a multiset decrease.
\end{hblemma}

\begin{hbtheorem}{thm:p07-ch88-check-domain-rank}{proved=annex:p7.ch88.check-domain-rank}{A typed boundary for the printed domain-rank measure}
In the valid open environment \(A:\ast,a:A,C:\ast\), put
\[
 L=(\lambda x:A)(\lambda B:\ast)B,\qquad
 M=(La)C,\qquad M'=((\lambda B:\ast)B)C.
\]
Both \(M\) and \(M'\) have source type \(\ast\), even in the
corrected restricted variant. An innermost permitted logical contraction
gives \(\Gamma\vdash M\triangleright M'\).
The printed multiset of domain ranks changes from \(\{0\}\) to
\(\{1\}\). Consequently the source's domain-rank multiset is not a
decreasing measure for these rules. The already proved normalization
theorem is independent of that measure.
\end{hbtheorem}

\begin{hblemma}{thm:p07-ch88-check-raw-reflection}{proved=annex:p7.ch88.check-raw-reflection}{Raw beta reflection and unique normal forms}
A target compatible beta step starting at an encoded scoped source term
reflects to exactly one scoped source beta step. Every target reduction
path from that image remains in the raw source image. Therefore the full
compatible raw source relation is confluent. On any source image supplied
with a CC typing certificate, full normalization terminates and has a
unique indexed normal form. No inverse of source typing or typed equality
is part of this assertion.
\end{hblemma}

\begin{hbdefinition}{def:p07-ch88-check-normal-variant}{The exact no-conversion normal-type variant}
Write \(\Gamma\vdash_n\Delta\) and \(\Gamma\vdash_n M:P\) for
the variant on Sheet~\ref{sheet:p07-ch88-normal-types}. It has the
eight p.~112 construction rules, normalized declaration annotations and
normalized application result types; it has no typing-conversion or
typed-equality rule. The notation \(\mathcal N(E)\) is the result of
a fixed full compatible beta-normalizing strategy, when that strategy
terminates, not a total function on arbitrary raw terms. The theorem
below proves that its calls at derived rule premises terminate.

Call a source expression \emph{annotation-normal} if every product or
lambda domain annotation is in full beta normal form, and its other
children are annotation-normal recursively. This is a finite syntactic
test; it does not normalize a possibly divergent input. Every derivable
environment annotation in this variant is fully normal.
\end{hbdefinition}

\begin{hbtheorem}{thm:p07-ch88-check-normal-transport}{proved=annex:p7.ch88.check-normal-transport}{All-rule transport of the normal-type variant}
Every normal-variant validity maps to \(T\Delta:\Box\), and every
normal-variant typing maps to \(TM:TP\), in exact two-sort CC and its
valid encoded context. All rule calls to \(\mathcal N\) are defined.
The subject is annotation-normal, the environment has normal annotations,
and every resulting source type is fully normal.
The target is exactly \(\ast:\Box\) with its four declared product
triples, not a universe extension or a proof-erased calculus.
\end{hbtheorem}

\begin{hbdefinition}{def:p07-ch88-check-algorithm}{A total algorithm for the normal-type variant}
Reject malformed or out-of-scope input before doing any normalization.
Validate a declaration list oldest-first, retaining the derivation and
CC certificate of each accepted prefix. In a checked prefix use these
mutually structural operations:
\begin{enumerate}
\item \(\mathrm{Sort}(A)\) first requires that \(A\) is syntactically
      fully normal. For a logical context it calls \(\mathrm{Valid}(A)\).
      For an object it calls \(\mathrm{Infer}(A)\) and requires the
      returned normal type to be literally \(\ast\).
\item \(\mathrm{Valid}(\ast)\) succeeds. For \([x:A]\Delta\),
      check \(\mathrm{Sort}(A)\), append its checked declaration,
      and check \(\mathrm{Valid}(\Delta)\) in that prefix.
\item \(\mathrm{Infer}(k)\) returns the checked context lookup.
      For an object product \([x:A]B\), check \(\mathrm{Sort}(A)\),
      infer \(B\) in the checked extension, and require its type
      literally \(\ast\); return \(\ast\).
\item For \((\lambda x:A)M\), check \(\mathrm{Sort}(A)\), infer
      the body type \(P\) in the checked extension, and return
      \([x:A]P\).
\item For \((MN)\), infer the normal function type. It must be a
      literal product \([x:A]B\). Infer the normal argument type
      and require literal equality with \(A\). Return
      \(\mathcal N([N/x]B)\). This last call is made only after the
      preceding successes have produced its target sort certificate.
\end{enumerate}
Any failed structural or equality check returns rejection. Each
recursive syntax call visits a proper child of its current original
expression, without rechecking the whole environment. Context validation
is a finite prefix loop. Lookup and comparison use the corrected indexed
operations, so equality of outputs is literal indexed equality, not
agreement of arbitrarily chosen binder names.
\end{hbdefinition}

\begin{hbtheorem}{thm:p07-ch88-check-normal-complete}{proved=annex:p7.ch88.check-normal-complete}{Soundness, completeness and termination of that checker}
The algorithm is total on finite scoped inputs and decides exactly the
normal variant's validity and typing judgments. On success it constructs
an actual derivation and, for an object, its unique literal normal type.
Thus supplied-type checking is equality with that returned type.

More precisely, in a checked annotation-normal environment, if
\(\widehat\Gamma\vdash_{\rm CC}TM:U\) for an annotation-normal
source object \(M\), the algorithm succeeds with a type \(P\) and
\(TP=_\beta U\). A source logical context typed in CC at \(\Box\)
is accepted by \(\mathrm{Valid}\). This is proved by structural
generation in CC, including the application product-head and normalized
domain cases; it is not an assumed inverse to the full source bridge.
Any encoded source environment which is CC-valid and has normal
annotations is likewise accepted by context validation.

Every derived \(\Gamma\vdash_n M:P\) has the actual source type
certificate \(\Gamma\vdash_n P:\ast\), if \(P\) is an object,
or \(\Gamma\vdash_n P\), if it is a logical context. It also has
\(\Gamma\vdash_n\mathcal N(M):P\). Preservation here is to the
whole normal form, not to every intermediate raw beta contractum.
\end{hbtheorem}

\begin{hbtheorem}{thm:p07-ch88-check-normal-step}{proved=annex:p7.ch88.check-normal-step}{Why normal types are a separate presentation}
In the valid normal environment \(A:\ast,a:A\), let
\[
 K=(\lambda x:A)A,\qquad
 V=((\lambda F:[x:A]\ast)(\lambda y:Fa)y)K.
\]
The normal variant derives \(V:[y:A]A\). Its raw root contractum
\((\lambda y:Ka)y\) has a nonnormal annotation and is not typable
in that variant. Its whole normal form \((\lambda y:A)y\) is typable.
Thus this variant is not raw-beta subject-preserving, and replacing a
raw full-calculus judgment by an unchanged normal-variant judgment is
not an adequacy theorem.
\end{hbtheorem}

\begin{proofboundary}{def:p07-ch88-check-full-boundary}{The exact ordinary-full checking dependency still open here}
For an annotation-normal full-source environment and subject, the proved
forward full transport followed by the new checker gives a normal-variant
derivation at the normal form of its old type. The converse ordinary-full
judgment and the checker on arbitrary full-source annotations still need
ordinary-full source type correctness, type preservation and completeness of typed source
equality on beta-convertible sorted endpoints. Raw image reflection alone
does not supply those derivations. They are not assumed here, and no
complete ordinary-full or restricted typed-equality checker is claimed.
The proved complete algorithm belongs to the explicitly declared p.~112
no-conversion variant.
\end{proofboundary}
